\subsection{A Slater point and a bound on the cap multipliers}\label{sec:P01}
From Lemma~\ref{lem:maximizador} we use only the existence and uniqueness of the canonical maximizer.
We neither assume that it lies in the interior nor assert that its weights coincide with the square
occupancy probabilities in a solution chosen uniformly at random.
Recall that $\mathcal D$ is the union of the two diagonal families of the board,
with $|\mathcal D|=4n-2$, and that
\[
 m_d(p)=\sum_{c\in d}p_c,\qquad h_d(p)=1-m_d(p).
\]
The polytope $K_n$ and the functional $H_n$ are those of~\eqref{eq:Kn} and~\eqref{eq:Hn}, with the
convention
$0\log0=0$. We denote by $p=\pi^{(n)}$ the unique maximizer of $H_n$ on $K_n$.

\paragraph{A Slater point for $n\ge32$.}
Let $r_i=(i-1/2)/n$ and
\[
 q_{ij}=\frac1n\left[1-\frac15\cos(2\pi r_i)\cos(2\pi r_j)\right].
\]
The sum of the cosines at the $n$ midpoints is zero for $n\ge2$.
Thus the row and column sums of $q$ are exactly one, and
$4/(5n)\le q_{ij}\le6/(5n)$.
The strictly admissible density $\nu^\sharp(r,c)=1-\frac15\cos(2\pi r)\cos(2\pi c)$
from~\cite[Remark~2.4]{Perfil} has, for $0\le d\le1$, the marginal
\[
 m_D[\nu^\sharp](d)=1-d-\frac15\left[
       \frac{1-d}{2}\cos(2\pi d)-\frac{\sin(2\pi d)}{4\pi}\right].
\]
Its derivative is at most $-1+(1+\pi)/5<0$, so $m_D[\nu^\sharp]\le9/10$.
By symmetry, the same bound holds throughout both families.
Each diagonal sum of $q$ is a midpoint quadrature for
$\nu^\sharp$ on the corresponding fiber, whose endpoints are multiples of $1/n$.
The parameter length of the fiber is at most one, and the Lipschitz constant along it
is at most $4\pi/5$. The quadrature error is thus
at most $(4\pi/5)/(2n)\le1/20$ for $n\ge32$, using $\pi<4$.
Consequently, with $M_0=6/5$ and $\kappa=1/20$,
\[
 q_{ij}\le M_0/n<1/8,\qquad h_d(q)\ge\kappa.       \tag{P1.1}
\]
This gives a Slater point in the affine space defined by the row and column sums.

\paragraph{Positivity in the presence of the cap.}
Take $p(t)=(1-t)p+tq$. If $p_c=0$, its spatial entropy contributes
$-q_c t\log t+O(t)$ to $H_n(p(t))-H_n(p)$. If $h_d(p)=0$, its
diagonal entropy contributes $-h_d(q)t\log t+O(t)$. The other entropy terms
contribute $O(t)$, as does $3(1-p)\log(1-p)$, whose derivative
is bounded on $[0,1/4]$. If any of these equalities holds, the
difference is $-At\log t+O(t)>0$ for small $t>0$, with $A>0$,
contradicting maximality. Hence
\[
 p_{ij}>0,\qquad h_d(p)>0.                         \tag{P1.2}
\]

\paragraph{Optimality conditions with cap multipliers.}
By (P1.2), the functional is differentiable near $p$, and the only potentially
active inequalities are $p_{ij}\le1/4$.
Slater's condition gives finite multipliers $\alpha_i,\beta_j$ and
$\mu_{ij}\ge0$. One of the row and column equalities is redundant and is omitted
when applying the Karush--Kuhn--Tucker (KKT) conditions~\cite[Ch.~5]{BV}. If $D(i,j),A(i,j)$ are the
two diagonals through the square, we obtain
\[
 -\log p_{ij}-3\log(1-p_{ij})+\log h_{D(i,j)}+\log h_{A(i,j)}-2
 =\alpha_i+\beta_j+\mu_{ij},\quad
 \mu_{ij}(1/4-p_{ij})=0.                          \tag{P1.3}
\]
Let $x_{ij}=np_{ij}$, $\ell_d=-\log h_d(p)\ge0$, and
$d_0=3\log(4/3)$. Multiplying (P1.3) by $q_{ij}-p_{ij}$ and
summing cancels the terms from the equality constraints and the constant $-2$. This gives
\[
 0=E+F+Q+U,                                      \tag{P1.4}
\]
\[
 E=\sum_{ij}(p_{ij}-q_{ij})\log x_{ij},\qquad
 F=3\sum_{ij}(p_{ij}-q_{ij})\log(1-p_{ij}),
\]
\[
 Q=\sum_{d\in\mathcal D}(h_d(q)-h_d(p))\ell_d,\qquad
 U=\sum_{ij}\mu_{ij}(1/4-q_{ij}).
\]
In particular, $U$ enters with a plus sign: by complementarity, the derivative in the direction $q-p$
is $-U$.

\paragraph{A global bound including the cap multipliers.}
Convexity gives $\sum p_{ij}\log(np_{ij})\ge0$. Moreover,
\[
 \sum q_{ij}\log x_{ij}\le\frac{M_0}{n}\sum\log^+x_{ij}
 \le\frac{M_0}{n}\sum x_{ij}=M_0n,
\]
so $E\ge-M_0n$. Since $\log(1-p)\le0$,
\[
 F\ge3\sum p_{ij}\log(1-p_{ij})\ge-d_0n.
\]
The inequality $h(-\log h)\le1/e$ and (P1.1) give
\[
 Q\ge\kappa\sum_d\ell_d-4n/e,\qquad
 U\ge\frac18\sum_{ij}\mu_{ij}.
\]
By (P1.4), setting $C_0=M_0+d_0+4/e$ yields
\[
 \kappa\sum_d\ell_d+\frac18\sum_{ij}\mu_{ij}\le C_0n.
\]
Since $\kappa\le1/8$, setting $T=C_0/\kappa$ gives
\[
 \frac1n\left(\sum_d\ell_d+\sum_{ij}\mu_{ij}\right)\le T.    \tag{P1.5}
\]
This estimate does not assume that any multiplier $\mu$ vanishes.
